\paragraph{Cayley 对数势、Poisson 原函数与有限缺陷扫描迹的统一闭式链}

\begin{theorem}[Cayley 模长的 Poisson--尺度积分表示与 Fourier 指纹]\label{thm:xi-cayley-modulus-poisson-fourier-fingerprint}
\leanverified{paper\_xi\_cayley\_modulus\_poisson\_fourier\_fingerprint}
设 Cayley 紧化
\[
\mathcal C(t):=\frac{1+\mathrm{i}t}{1-\mathrm{i}t},\qquad
\mathcal C^{-1}(w)=-\mathrm{i}\,\frac{w-1}{w+1},
\]
并定义 Poisson--Cauchy 核
\[
\mathsf P_s(x):=\frac{1}{\pi}\frac{s}{x^2+s^2}\qquad(s>0,\ x\in\RR).
\]
对任意 \(\delta\in(0,1)\) 与 \(x\in\RR\)，令
\[
u_\delta(x):=-\log\bigl|\mathcal C(x+\mathrm{i}\delta)\bigr|.
\]
则成立：
\begin{enumerate}
  \item
  \[
  u_\delta(x)=\frac12\log\frac{x^2+(1+\delta)^2}{x^2+(1-\delta)^2};
  \]
  \item
  \[
  u_\delta(x)=\pi\int_{1-\delta}^{1+\delta}\mathsf P_s(x)\,\dd s;
  \]
  \item 取 Fourier 约定 \(\widehat f(\xi)=\int_{\RR}e^{-\mathrm{i}\xi x}f(x)\,\dd x\)，则对 \(\xi\neq 0\),
  \[
  \widehat{u_\delta}(\xi)=
  \pi\frac{e^{-(1-\delta)|\xi|}-e^{-(1+\delta)|\xi|}}{|\xi|},
  \]
  并且
  \[
  \widehat{u_\delta}(0)=\int_{\RR}u_\delta(x)\,\dd x=2\pi\delta
  \]
  （连续延拓意义下）。
\end{enumerate}
\end{theorem}
\begin{proof}
对 \(t=x+\mathrm{i}\delta\) 直接计算
\[
\bigl|\mathcal C(t)\bigr|^2
=\frac{|1+\mathrm{i}t|^2}{|1-\mathrm{i}t|^2}
=\frac{x^2+(1-\delta)^2}{x^2+(1+\delta)^2},
\]
取负对数即得第一式。

再由
\[
\int_a^b\mathsf P_s(x)\,\dd s
=\frac{1}{\pi}\int_a^b\frac{s}{x^2+s^2}\,\dd s
=\frac{1}{2\pi}\log\frac{x^2+b^2}{x^2+a^2},
\]
取 \(a=1-\delta,\ b=1+\delta\) 即得第二式。

最后利用 \(\widehat{\mathsf P_s}(\xi)=e^{-s|\xi|}\)，对第二式逐点 Fourier 变换，得到
\[
\widehat{u_\delta}(\xi)
=\pi\int_{1-\delta}^{1+\delta}e^{-s|\xi|}\,\dd s
=\pi\frac{e^{-(1-\delta)|\xi|}-e^{-(1+\delta)|\xi|}}{|\xi|}\quad(\xi\neq 0).
\]
令 \(\xi\to 0\) 得连续延拓值 \(2\pi\delta\)，即 \(\widehat{u_\delta}(0)=\int u_\delta\)。
\end{proof}

\begin{corollary}[深度扫描面积律与低频迹公式]\label{cor:xi-cayley-scan-area-law-lowfreq-trace}
\leanverified{paper\_xi\_cayley\_scan\_area\_law\_lowfreq\_trace}
在定理 \ref{thm:xi-cayley-modulus-poisson-fourier-fingerprint} 记号下，设有限离散缺陷谱
\[
\nu=\sum_{j=1}^{\kappa}m_j\,\delta_{(\gamma_j,\delta_j)},
\qquad
m_j\in\NN,\ \delta_j\in(0,1),\ \gamma_j\in\RR,
\]
并定义 Jensen--Cayley 扫描势
\[
J_\nu(x):=\sum_{j=1}^{\kappa}m_j\,u_{\delta_j}(x-\gamma_j).
\]
则
\[
\int_{\RR}J_\nu(x)\,\dd x
=2\pi\sum_{j=1}^{\kappa}m_j\delta_j,
\]
且
\[
\widehat{J_\nu}(0)=2\pi\sum_{j=1}^{\kappa}m_j\delta_j,
\qquad
\lim_{\xi\to 0}\widehat{J_\nu}(\xi)=2\pi\sum_{j=1}^{\kappa}m_j\delta_j.
\]
\end{corollary}
\begin{proof}
对每一项 \(u_{\delta_j}(x-\gamma_j)\) 使用定理
\ref{thm:xi-cayley-modulus-poisson-fourier-fingerprint} 的 Poisson 原函数表示，并利用
\(\int_{\RR}\mathsf P_s(x-\gamma_j)\dd x=1\)，交换积分与有限求和即得面积律。
再由 \(\widehat{J_\nu}(0)=\int J_\nu\) 得低频迹公式。
\end{proof}

\begin{theorem}[有限缺陷口径下的边界 Jensen 潜势极限与显式分解]\label{thm:xi-jensen-boundary-potential-finite-defect-explicit}
沿用 \eqref{eq:xi-comoving-shifted-jensen-defect} 的 \(D_{\zeta,x_0}(\varrho)\) 记号。
设 \(\Xi_\zeta\) 在离临界线侧仅有有限个零点（按重数计）
\[
\rho_j=\beta_j+\mathrm{i}\gamma_j,\qquad
\beta_j<\frac12,\qquad
\delta_j:=\frac12-\beta_j\in(0,1),\qquad
m_j\in\NN\quad(1\le j\le \kappa).
\]
对固定 \(x_0\in\RR\)，令
\[
w_j(x_0):=\mathcal C\bigl((\gamma_j-x_0)+\mathrm{i}\delta_j\bigr)\in\mathbb D.
\]
则极限
\[
J_\zeta(x_0):=\lim_{\varrho\uparrow 1}D_{\zeta,x_0}(\varrho)
\]
存在，并且
\[
J_\zeta(x_0)
=\sum_{j=1}^{\kappa}m_j\Bigl(-\log|w_j(x_0)|\Bigr)
=\sum_{j=1}^{\kappa}m_j\,u_{\delta_j}(x_0-\gamma_j).
\]
此外 \(J_\zeta\in L^1(\RR)\cap C^\omega(\RR)\)，并有尾部展开
\[
J_\zeta(x_0)
=2\sum_{j=1}^{\kappa}\frac{m_j\delta_j}{(x_0-\gamma_j)^2}
+O(|x_0|^{-4})
=\frac{2}{x_0^2}\sum_{j=1}^{\kappa}m_j\delta_j+O(|x_0|^{-3})
\quad(|x_0|\to\infty).
\]
\end{theorem}
\begin{proof}
由定理 \ref{thm:xi-comoving-jensen-defect-kernel-decomposition}，
\[
D_{\zeta,x_0}(\varrho)=\sum_{j=1}^{\kappa}m_j\Bigl[\log\frac{\varrho}{|w_j(x_0)|}\Bigr]_+.
\]
因 \(0<|w_j(x_0)|<1\) 且和式有限，故当 \(\varrho\) 充分接近 \(1\) 时，所有项均处于“激活”状态，且
\[
\Bigl[\log\frac{\varrho}{|w_j(x_0)|}\Bigr]_+\to -\log|w_j(x_0)|.
\]
于是得到极限存在与显式分解。

每一项 \(u_{\delta_j}(x_0-\gamma_j)\) 在实轴上实解析，且按定理
\ref{thm:xi-cayley-modulus-poisson-fourier-fingerprint} 属于 \(L^1\)，有限求和得 \(J_\zeta\in L^1\cap C^\omega\)。
最后由
\[
u_\delta(y)=\frac12\log\frac{y^2+(1+\delta)^2}{y^2+(1-\delta)^2}
=\frac{2\delta}{y^2}+O(|y|^{-4})
\]
逐项代入并求和即得尾部展开。
\end{proof}

\begin{theorem}[Jensen 潜势的 Fourier--Laplace 指纹与有限缺陷注入性]\label{thm:xi-jensen-fourier-laplace-fingerprint-injectivity}
\leanverified{paper\_xi\_jensen\_fourier\_laplace\_fingerprint\_injectivity}
在定理 \ref{thm:xi-jensen-boundary-potential-finite-defect-explicit} 的设定下，对 \(\xi\neq 0\) 有
\[
\widehat{J_\zeta}(\xi)
=\sum_{j=1}^{\kappa}m_j\,
\pi\frac{e^{-(1-\delta_j)|\xi|}-e^{-(1+\delta_j)|\xi|}}{|\xi|}
e^{-\mathrm{i}\gamma_j\xi}.
\]
令 \(s>0\) 并定义
\[
G(s):=\frac{s}{\pi}\,\widehat{J_\zeta}(s)
=\sum_{j=1}^{\kappa}m_j
\Bigl(e^{-(1-\delta_j+\mathrm{i}\gamma_j)s}
-e^{-(1+\delta_j+\mathrm{i}\gamma_j)s}\Bigr).
\]
则 \(G\) 是有限指数和；其指数集
\[
\Lambda=\{\,1-\delta_j+\mathrm{i}\gamma_j,\ 1+\delta_j+\mathrm{i}\gamma_j:\ 1\le j\le\kappa\,\}
\]
及对应系数由 \(G\) 唯一确定。因而映射
\[
\{(\gamma_j,\delta_j,m_j)\}_{j=1}^{\kappa}\longmapsto J_\zeta(\cdot)
\]
在有限缺陷类别上为注入。
\end{theorem}
\begin{proof}
第一式由定理 \ref{thm:xi-jensen-boundary-potential-finite-defect-explicit} 的分解与定理
\ref{thm:xi-cayley-modulus-poisson-fourier-fingerprint} 的 Fourier 闭式逐项求和得到。

若两组有限缺陷数据给出同一 \(J_\zeta\)，则对应差函数满足 \(G(s)\equiv 0\)（\(s>0\)），可写为
\[
\sum_{\ell=1}^{N}c_\ell e^{-\lambda_\ell s}=0,\qquad
\Re\lambda_\ell>0,\ \lambda_\ell\ \text{互异}.
\]
有限指数函数组在开区间上线性无关，故全部 \(c_\ell=0\)。
因此指数与系数逐项一致，从而 \((\gamma_j,\delta_j,m_j)\) 的多重集唯一确定。
\end{proof}

\begin{proposition}[Jensen--Poisson 互织算子：分数导数产生双尺度 Poisson 剖面]\label{prop:xi-jensen-poisson-intertwiner-fractional-derivative}
仍在定理 \ref{thm:xi-jensen-boundary-potential-finite-defect-explicit} 设定下，令 \(|D|\) 为 Fourier 乘子 \(|\xi|\)。
则在温和分布意义下有
\[
|D|\,J_\zeta
=\pi\sum_{j=1}^{\kappa}m_j\Bigl(
\mathsf P_{1-\delta_j}(\cdot-\gamma_j)-\mathsf P_{1+\delta_j}(\cdot-\gamma_j)\Bigr).
\]
进一步，对任意 \(t>0\)（其中 \(e^{-t|D|}f=\mathsf P_t*f\)）有
\[
\mathsf P_t*|D|J_\zeta
=\pi\sum_{j=1}^{\kappa}m_j\Bigl(
\mathsf P_{t+1-\delta_j}(\cdot-\gamma_j)-\mathsf P_{t+1+\delta_j}(\cdot-\gamma_j)\Bigr).
\]
\end{proposition}
\begin{proof}
由定理 \ref{thm:xi-cayley-modulus-poisson-fourier-fingerprint}，
\[
\widehat{u_\delta}(\xi)=
\pi\frac{e^{-(1-\delta)|\xi|}-e^{-(1+\delta)|\xi|}}{|\xi|}.
\]
两侧乘以 \(|\xi|\) 即得
\[
\widehat{|D|u_\delta}(\xi)=\pi\bigl(e^{-(1-\delta)|\xi|}-e^{-(1+\delta)|\xi|}\bigr)
=\pi\bigl(\widehat{\mathsf P_{1-\delta}}-\widehat{\mathsf P_{1+\delta}}\bigr)(\xi).
\]
反演并对有限和线性叠加，得到第一式。
再乘以 \(e^{-t|\xi|}\)（等价于卷积 \(\mathsf P_t\)）即得第二式。
\end{proof}

\begin{theorem}[有限缺陷口径下的 RH 等价判据：Jensen 扫描迹公式]\label{thm:xi-finite-defect-rh-criterion-jensen-scan-trace}
在定理 \ref{thm:xi-jensen-boundary-potential-finite-defect-explicit} 的设定下，定义
\[
\mathcal J(\varrho):=\frac{1}{2\pi}\int_{\RR}D_{\zeta,x_0}(\varrho)\,\dd x_0\qquad(0<\varrho<1).
\]
则极限存在且
\[
\lim_{\varrho\uparrow 1}\mathcal J(\varrho)=\sum_{j=1}^{\kappa}m_j\delta_j.
\]
特别地，在该有限缺陷口径下，下述命题等价：
\begin{enumerate}
  \item Riemann 假设成立；
  \item \(\displaystyle \lim_{\varrho\uparrow 1}\mathcal J(\varrho)=0\)；
  \item \(\displaystyle \lim_{\xi\to 0}\widehat{J_\zeta}(\xi)=0\)。
\end{enumerate}
\end{theorem}
\begin{proof}
由定理 \ref{thm:xi-jensen-boundary-potential-finite-defect-explicit}，
\(D_{\zeta,x_0}(\varrho)\to J_\zeta(x_0)\)；
并由有限分解与 \(u_{\delta_j}\in L^1\) 得支配收敛可用，从而
\[
\lim_{\varrho\uparrow 1}\mathcal J(\varrho)
=\frac{1}{2\pi}\int_{\RR}J_\zeta(x_0)\,\dd x_0.
\]
再由推论 \ref{cor:xi-cayley-scan-area-law-lowfreq-trace} 对 \(J_\zeta=\sum_j m_j u_{\delta_j}(\cdot-\gamma_j)\) 的面积律可得
\[
\frac{1}{2\pi}\int_{\RR}J_\zeta=\sum_{j=1}^{\kappa}m_j\delta_j.
\]
于是极限公式成立。

若 RH 成立，则不存在 \(\delta_j>0\)，故和为 \(0\)；反之若该和为 \(0\)，因 \(m_j,\delta_j\) 皆非负且 \(\delta_j>0\) 表示离线偏移，故只能全部不存在，从而 RH 成立。
第三条由 \(\widehat{J_\zeta}(0)=\int J_\zeta\) 与连续性直接等价于第二条。
\end{proof}

\begin{proposition}[有限半径扫描迹闭式：高度消去与深度谱锁定]\label{prop:xi-finite-rho-scan-trace-closed-form-depth-only}
\leanverified{paper\_xi\_finite\_rho\_scan\_trace\_closed\_form\_depth\_only}
令 \(\delta\in(0,1)\)、\(\varrho\in(0,1)\)，并定义截断核
\[
K_{\varrho,\delta}(x):=\Bigl[\log\frac{\varrho}{|\mathcal C(x+\mathrm{i}\delta)|}\Bigr]_+.
\]
则有：
\begin{enumerate}
  \item 若 \(\varrho\le \frac{1-\delta}{1+\delta}\)，则 \(K_{\varrho,\delta}\equiv 0\)；
  \item 若 \(\varrho> \frac{1-\delta}{1+\delta}\)，则存在唯一 \(X_\star(\varrho,\delta)>0\) 使
  \[
  \operatorname{supp}K_{\varrho,\delta}=[-X_\star,X_\star],\qquad
  X_\star^2=\frac{\varrho^2(1+\delta)^2-(1-\delta)^2}{1-\varrho^2},
  \]
  且
  \begin{align*}
  \int_{\RR}K_{\varrho,\delta}(x)\,\dd x
  =&\ 2X_\star\log\varrho
  +X_\star\log\frac{X_\star^2+(1+\delta)^2}{X_\star^2+(1-\delta)^2}\\
  &\ +2\Bigl[(1+\delta)\arctan\!\frac{X_\star}{1+\delta}
  -(1-\delta)\arctan\!\frac{X_\star}{1-\delta}\Bigr].
  \end{align*}
\end{enumerate}
进一步，在有限缺陷口径下，
\[
\int_{\RR}D_{\zeta,x_0}(\varrho)\,\dd x_0
=\sum_{j=1}^{\kappa}m_j\int_{\RR}K_{\varrho,\delta_j}(x)\,\dd x,
\]
其右端与全部 \(\gamma_j\) 无关：任意固定 \(\varrho\) 的 \(x_0\)-扫描迹仅依赖深度谱 \(\{(\delta_j,m_j)\}\)。
\end{proposition}
\begin{proof}
由
\[
\bigl|\mathcal C(x+\mathrm{i}\delta)\bigr|^2
=\frac{x^2+(1-\delta)^2}{x^2+(1+\delta)^2},
\]
不等式 \(|\mathcal C(x+\mathrm{i}\delta)|<\varrho\) 等价于
\[
(1-\varrho^2)x^2<\varrho^2(1+\delta)^2-(1-\delta)^2.
\]
故得到阈值与支集公式。

在支集内，
\[
K_{\varrho,\delta}(x)
=\log\varrho+\frac12\log\bigl(x^2+(1+\delta)^2\bigr)-\frac12\log\bigl(x^2+(1-\delta)^2\bigr).
\]
使用原函数
\[
\int_0^X\log(x^2+a^2)\,\dd x
=X\log(X^2+a^2)-2X+2a\arctan(X/a),
\]
在 \([0,X_\star]\) 上计算并利用偶对称乘以 \(2\)，得到闭式积分表达。

最后一式来自定理 \ref{thm:xi-comoving-jensen-defect-kernel-decomposition} 的逐项可加性与平移不变性：每个核仅以 \(x=\gamma_j-x_0\) 进入，对 \(x_0\) 全轴积分后 \(\gamma_j\) 自动消去。
\end{proof}
